\chapter{把一张取数表，逐行翻译成矩阵}\label{app:upper}
\section{上半带的大载体从哪里来}
设 $n=3p+s$，$r=p+s$。核心是 $(A^p,B^p,C^r)$。现在载体有 $r>p$ 份，如果想全部换成一个新均值，做一次平均只会产生 $p$ 份，显然不够。因此先做两组\emph{组成相同、位置不交}的平均，共产生 $2p$ 份新值 $D$。

每组从 $A,B,C$ 取 $i,j,k$ 份，$i+j+k=p$。新值为
\[D=\frac{iA+jB+kC}{p}
=\frac{(i-j)a+(rj-pk)z}{p}.
\]
保留 $r$ 份 $D$ 作新载体，其余恰有 $2p$ 个位置，分成两组输出即可。这件事的全部自由度，就是来源三项计数，以及第一输出组的四项计数。

\section{先把库存列出来，再做任何代数}
两组来源用去 $2i,2j,2k$ 份旧值。保留 $r$ 份 $D$ 后，剩余四类库存是
\[(p-2i, p-2j, r-2k, 2p-r).
\]
第一输出取 $(x,y,u,v)$。它必须逐项不超过库存，且四项和为 $p$。第二输出自动是剩下的 $p$ 项。此处 $y$ 是一项计数，暂时不是先前的见证坐标 $a+pz$；读局部公式时要按本段定义使用。

新载体满足 $-pz'=D$，所以矩阵第二行直接得到。第一输出的均值是 $a'$，所以把 $(xA+yB+uC+vD)/p$ 展开，就得到第一行。
\InsertProof{lem:count-interface}

\section{一整套数字：十一平均、三十四个位置}
取 $p=11,s=1,r=12,n=34$。来源 $(i,j,k)=(4,3,4)$。做两次相同组成的平均后，留下十二份 $D$，库存为
\[(3,5,4,10).
\]
可选择两个不同输出方案 $(3,4,1,3)$ 和 $(2,3,3,3)$。每个方案各自形成一条四次平均的完整返回。

$\alpha=i-j=1$，$\beta=rj-pk=-8$。矩阵为
\[
J_f=\begin{pmatrix}-8&383\\-1&8\end{pmatrix},qquad
J_g=\begin{pmatrix}-8&9\\-1&8\end{pmatrix}.
\]
平方分别为 $-319I$ 和 $55I$，故 $f=-29,g=5$，且 $g-f=34=\alpha n$。它们的商给倍率 $-5/29$，环为 $\Z[1/145]$。再算
\[H=n(2j+\alpha)-2p^2=34\cdot7-242=-4,
\]
得到上剪切 $U(-136/5)$。横向系数 $\alpha=1$，根步长是 $4nR$，便能进入共同终端允许的偶数层。

这张表里的数字不是“某几个输入能平均”的例子。它给出的矩阵等式对任意参数 $a,z$ 恒成立，所以这一个计数表可以处理整个无界的二维配置族。

\section{两张很像的表，为什么比一张强很多}
从第一输出改成 $(x-1,y-1,u+2,v)$，总项数没变，$x-y$ 也没变，迹零条件仍成立；只有右上项少了 $pr+2p^2=pn$。

于是比较两反射会产生三角矩阵；再用换块的对称性消去三角矩阵的倍率，剩下一个剪切。若有几组表，利用剪切参数相加和裴蜀等式，就把它们的步长合成最大公因数。

下面的推论把后面的工作归结为两个清楚的目标：横向系数的最大公因数为一，根系数的最大公因数足够小。满足这两项，群论和终端部分已经不用重新证明。
\InsertProof{cor:resource}
